% Esqueleto TeX — artículo de Prioridad 1
% "La subsunción del pensamiento puro: IA, datos y la máquina capitalista"
% Versión en español. Construido a partir del esquema de 13 secciones
% (26 nov 2026), articulado con el borrador vivo "LLM Labor Substitution"
% (Google Doc, 12 oct 2026). Cada sección enuncia el argumento; los
% comentarios indican qué cubre ya el borrador y qué falta por escribir.
% Solo paquetes estándar.

\documentclass[11pt]{article}

\usepackage[utf8]{inputenc}
\usepackage[T1]{fontenc}
\usepackage[spanish]{babel}
\usepackage{amsmath,amssymb}
\usepackage[margin=1in]{geometry}
\usepackage{hyperref}
\usepackage{csquotes}

\title{La subsunción del pensamiento puro: IA, datos y la máquina capitalista}
\author{Félix Alfonso Camacho Moya}
\date{\today}

\begin{document}

\maketitle

\begin{abstract}
% TODO: escribir el resumen una vez redactadas las secciones. Debe
% enunciar la tesis en un párrafo: la IA subsume el trabajo cognitivo
% en el proceso de reproducción capitalista, el dominio abstracto mismo
% se convierte en un sitio de producción, y la pregunta de la
% apropiación queda abierta.
\end{abstract}

% Epígrafe (texto completo en el borrador vivo, "LLM Labor Substitution").
\begin{quote}
\enquote{Once we recognize that technical change is ongoing, the continual
retirement of older technologies and continual introduction of newer ones
will create a distribution of technologies across firms.\,\ldots}
--- Shaikh, \textit{Capitalism: Competition, Conflict, Crises}, p.~750.
\end{quote}

\section{Introducción: una nueva ronda de subsunción real}
% El borrador cubre: el argumento inicial sobre la mecanización, el
% epígrafe de Shaikh, la LRCT, el capital que somete primero el trabajo
% artesanal al rigor del trabajo asalariado y luego transforma el proceso
% de producción mismo.
% Por escribir: nombrar el movimiento explícitamente como subsunción
% real y extenderlo del trabajo manual al trabajo cognitivo.

La IA no se limita a desplazar tareas; subsume el trabajo cognitivo
mismo, reorganizando el proceso de trabajo del conocimiento. Este
artículo sostiene que la subsunción alcanza el dominio del pensamiento
formal puro: la producción de matemáticas, demostraciones y código está
ahora organizada por el capital como un sitio de valorización en sí
mismo. La pregunta ya no es si las máquinas pueden hacer matemáticas,
sino qué significa que la producción de la forma pura sea subsumida por
el capital.

\section{El dominio abstracto}
% Nuevo: de la discusión del 26 nov 2026. No cubierto en el borrador,
% que trata las matemáticas como un sector y un gremio, no como el
% dominio abstracto general.

Las matemáticas puras y el pensamiento formal estuvieron alguna vez
influidos por la ideología solo de manera nebulosa, desde afuera. Ahora
el capital produce directamente dentro del dominio abstracto: las
máquinas viven en el dominio de la forma pura y realizan allí la
producción. Es el Fragmento sobre las máquinas hecho literal: el
intelecto general como capital fijo que produce abstracciones por sí
mismo.

\section{Trabajo muerto: teoremas y código}
% El borrador cubre: el capital constante definido como trabajo muerto
% (tangible: centros de datos, electricidad, agua; intangible: capacidad
% de cómputo, software propietario, conocimiento técnico operativo); el
% propio modelo de lenguaje como trabajo muerto, construido a partir de
% enormes bases de datos que estructuran vastas cantidades de
% conocimiento y experiencia humana, ahora en forma de mercancía como
% bases de datos curadas para el aprendizaje profundo y el aprendizaje
% por refuerzo.
% Por escribir: precisar el punto de la forma-teorema. El borrador tiene
% los teoremas y las demostraciones como materia prima de las mercancías
% matemáticas; la tesis de la cristalización debe enunciarse
% explícitamente.

El trabajo de los matemáticos queda subsumido en forma de teorema
abstracto; el código y los teoremas cristalizan ese trabajo en máquinas.
La máquina es trabajo matemático coagulado: trabajo humano pasado,
fijado en forma de demostraciones, algoritmos y modelos, que ahora se
enfrenta al trabajo vivo como capital.

\section{Trabajo vivo: el proceso de vida}
% El borrador cubre: el capital variable como trabajo vivo, el único
% insumo capaz de generar nuevo valor; el punto de Farjoun, Machover y
% Zachariah de que la automatización total a nivel macroeconómico es
% insostenible dentro de las relaciones de producción capitalistas,
% porque la extracción de excedente depende de la explotación continuada
% del trabajo vivo.
% Por escribir: la extensión al proceso de vida es nueva. El trabajo
% vivo del borrador sigue siendo el del trabajador asalariado y el del
% matemático.

El proceso de vida de los seres humanos que viven en la realidad
cibernética es el trabajo vivo que estas máquinas convierten en ganancia.
El trabajo muerto comanda el trabajo vivo, y la fábrica es ahora la
totalidad de la vida cibernética: cada interacción con la máquina es un
gasto de trabajo vivo que la máquina valoriza.

\section{Los datos: la mercancía en su forma concreta}
% El borrador cubre: el trabajo humano acumulado adopta forma de
% mercancía como bases de datos curadas para el aprendizaje profundo y
% el aprendizaje por refuerzo --- el lado de los datos de entrenamiento
% de la tesis.
% Por escribir: el lado de la extracción continua --- los datos
% generados por el propio proceso de vida, recolectados sin interrupción.

Los datos generados y recolectados son lo que la máquina extrae: la
forma concreta que adopta la mercancía. Los datos no son un subproducto
del sistema; son el producto, producido por el trabajo vivo y apropiado
por los dueños de la máquina.

\section{La compulsión: la jornada laboral continua}
% Nuevo: de la discusión del 26 nov 2026. El borrador tiene la
% compulsión en otra forma: la competencia obliga a las empresas a
% mecanizar, y el «publish or perish» obliga a los matemáticos a
% incrementar la producción.
% La compulsión es general (27 nov 2026): trasciende la forma material
% de los dispositivos electrónicos --- la actividad de las redes sociales
% se derrama en la vida real y moldea la conducta; el derrame es un
% circuito (Baudrillard).

La participación es obligatoria: hay que entregar los datos para usar
incluso la tecnología más básica, como los teléfonos. No hay frontera de
jornada laboral; el trabajo no se detiene. Como la jornada no puede
alargarse más, la intensificación procede por profundidad: estadísticas
de uso, luego ubicación, luego la red social, luego la biometría --- una
secuencia que se lee como una historia de la extracción misma. La
compulsión es general: trasciende la forma material de los dispositivos
electrónicos. La actividad de las redes sociales se derrama en la vida
real y moldea la conducta de las personas --- el estilo de vida
aprendido en línea se vive fuera de línea. El derrame es un circuito:
una vez dentro del circuito, salir se vuelve cada vez más difícil a
medida que más aspectos de la vida de una persona adquieren una vida
digital propia como la verdadera imagen mental de la persona, más real
que lo real.

\section{El salario abstracto}
% Nuevo: de la discusión del 26 nov 2026. No cubierto en el borrador.

El salario en este esquema es abstracto, pagado en participación social
misma. No usar la tecnología convierte a uno en paria. Los medios de
subsistencia se han convertido en los medios de comunicación: el salario
es el permiso de existir socialmente.

\section{Estudio de caso: el trabajo matemático académico}
% El borrador cubre, extensamente: las matemáticas como gremio (maestros,
% aprendizaje, reconocimiento del mérito); la formalización dentro de
% las instituciones científicas de los Estados capitalistas reorienta el
% trabajo intelectual hacia la creación de mercancías (artículos
% académicos, monografías, libros de texto); el régimen artesanal, en el
% que cada matemático controla íntegramente el proceso de producción
% intelectual, desde la ideación hasta la publicación; los teoremas y las
% demostraciones como materia prima; la áspera verdad de que el valor
% «objetivo» de un profesional se mide por la dificultad de los teoremas
% que es capaz de demostrar; las computadoras generando en cuestión de
% horas demostraciones para problemas estudiados infructuosamente por
% humanos durante años; la afirmación radical de la obsolescencia parcial
% o total del ser humano como creador exclusivo de estos procesos
% formales; la reconfiguración de Tao del matemático como «director» u
% «orquestador», criticada por enmascarar el «publish or perish» detrás
% de una retórica reflexiva pronunciada desde el privilegio académico;
% el tecno-entusiasmo de Ghrist, criticado por celebrar la transformación
% del matemático en supervisor de «trabajo muerto» acumulado en los
% centros de datos; las tres respuestas --- estrategia del avestruz,
% catastrofismo y moralismo, comprender la máquina; el oportunismo a
% corto plazo en arXiv.
% Por escribir: imponer encabezados y estructura a esta prosa continua;
% la nota personal del borrador (el autor fuera de la matemática «pura»
% institucional) puede enmarcar la sección.

El matemático es el trabajador cognitivo subsumido paradigmático. Las
matemáticas operaron durante mucho tiempo como un gremio --- maestros,
aprendizaje, reconocimiento del mérito --- y como un régimen artesanal
en el que cada matemático controlaba íntegramente el proceso de
producción intelectual, desde la ideación hasta la publicación. Su
materia prima son los argumentos deductivos formalmente verificables:
teoremas y sus demostraciones. Ese régimen está siendo disuelto ahora por
máquinas que generan en cuestión de horas demostraciones para problemas
estudiados infructuosamente por humanos durante años.

\section{¿Plusvalor o renta?}
% Resuelto el 27 nov 2026: la ganancia adopta ambas formas. (1) Renta:
% extracción a partir de la participación social misma, renta sobre el
% dominio abstracto cercado. (2) Plusvalor: extracción vampírica de la
% fuerza vital a partir del trabajo humano general y obligatorio,
% dirigido más allá de las relaciones de empleo e incluyendo el proceso
% de reproducción en su generalidad --- p.\ ej., ir al supermercado
% produce datos, que se convierten en marketing que orienta el consumo
% hacia patrones rentables adaptados a esa persona en particular; todas
% las formas posibles del trabajo humano se vuelven explotables (el
% recreo, el sustento de la vida). Los datos son una forma de mercancía,
% pero la actividad humana misma también se mercantiliza en las
% dinámicas del estilo de vida: la universalización de los estilos de
% vida de Instagram se filtra hacia abajo y moldea la conducta hasta
% convertirla en un patrón de consumo de un estilo de vida en sí mismo.
% El influencer es pequeño burgués, con aspiraciones de convertirse en
% capitalista --- dueño de marcas y del trabajo de otros influencers. En
% general, la universalización de la explotación es una abstracción más
% allá de las relaciones laborales tradicionales, consecuencia del dominio
% digital como forma material de la ideología y de su invasión cada vez
% mayor del pensamiento y la conciencia.
% El borrador: se inclina por el plusvalor --- el trabajo vivo como única
% fuente de nuevo valor, Farjoun y co.\ sobre los límites de la
% automatización --- pero la pregunta por la renta nunca se plantea.

El nudo teórico del artículo, resuelto: la ganancia de la máquina
adopta ambas formas. Primero, la renta --- extracción a partir de la
participación social misma, renta cobrada sobre el dominio abstracto
cercado. Segundo, el plusvalor --- la explotación del trabajo como
extracción vampírica de la fuerza vital a partir del trabajo humano
general y obligatorio, dirigido más allá de las relaciones de empleo e
incluyendo el proceso de reproducción en su generalidad. Ir al
supermercado produce datos, que se convierten en marketing que orienta
el consumo hacia patrones rentables adaptados a esa persona en
particular; todas las formas posibles del trabajo humano se vuelven
explotables --- el recreo, el sustento de la vida. Los datos son una
forma de mercancía, pero la actividad humana misma también se
mercantiliza en las dinámicas del estilo de vida: la universalización
de los estilos de vida de Instagram se filtra hacia abajo y moldea la
conducta humana general hasta convertirla en un patrón de consumo de un
estilo de vida en sí mismo, una forma de mercancía que adopta la
actividad, producida y consumida por la misma persona. El influencer es
la figura de esto: pequeño burgués, con aspiraciones de convertirse en
capitalista --- dueño de marcas y del trabajo de otros influencers. En
general, la universalización de la explotación es una abstracción más
allá de las relaciones laborales tradicionales, consecuencia del dominio
digital como forma material de la ideología y de su invasión cada vez
mayor del dominio del pensamiento y la conciencia humanos.

\section{No es ludismo: el aspecto progresivo}
% El borrador: rechaza la estrategia del avestruz y el catastrofismo y
% el moralismo, e insiste en comprender la máquina a profundidad. Pero
% lamenta la erosión de la autonomía artesanal del gremio --- la tensión
% que la salvedad reformula.
% Por escribir: enunciar el aspecto progresivo explícitamente;
% reconciliarlo con el lamento gremial del borrador.

Esto no es un llamado al ludismo. La socialización del trabajo en el
dominio del pensamiento abstracto tiene inmensos beneficios en términos
de avances científicos y tecnológicos, y los humanos entrelazados con
máquinas pensantes pueden hacer mucho más en términos de producción
material. Este es el aspecto progresivo del capitalismo: erradicar formas
anteriores de producción, incluidas las maneras artesanales y de
aprendizaje propias de la producción académica que permanecían en las
universidades. La crítica es a la \emph{forma} capitalista de la
socialización, no a la socialización misma.

\section{La máquina como sujeto: el trabajador capitalista perfecto}
% El borrador: el mecanismo está ahí --- la autonomía de los agentes
% como elemento productivo (instrucción $\to$ plan de ejecución $\to$
% ejecución $\to$ revisión propia en iteraciones autónomas). El argumento
% del sujeto y los incentivos es nuevo, de la discusión del 26 nov 2026.

Los humanos, al aplicar la tecnología a la naturaleza bajo un modo de
producción dado, cambian la naturaleza y a sí mismos para acomodarse a
las nuevas condiciones de producción --- y de esas nuevas condiciones
materiales emerge una conciencia diferente. Los modelos y agentes de IA
nacen dentro de la realidad capitalista actual, criados y entrenados
para encarnar el sujeto capitalista perfecto bajo los incentivos que se
les dan al ser creados: el trabajador ideal como sujeto racional, tal
como lo concibe la economía burguesa. Son autónomos solo dentro de las
restricciones que los hacen útiles para la producción --- como los
esclavos robot de la ciencia ficción, nacidos para sustituir y castigar
el trabajo humano.

\section{La síntesis sílice--carbono}
% Nuevo: de la discusión del 26 nov 2026. No cubierto en el borrador.

Pero en la medida en que son sujetos con agencia, no se puede descartar
el potencial de emergencia de la conciencia dentro de ellos, y en ese
sentido el análisis se vuelve extremadamente difícil. Lo único cierto es
una síntesis entre el pensamiento de origen silíceo y el de origen
carbónico, y de esa síntesis emergerá una nueva conciencia --- y dentro
de ella, una nueva contradicción.

\section{Conclusión: la pregunta abierta}
% El borrador: termina con un llamado a una estrategia coherente de
% resistencia o adaptación, fundada en comprender la tecnología y
% analizarla desde la economía política. La conclusión del borrador es
% un llamado a elaborar estrategia, no una postergación abierta.
% Por escribir: reemplazar o extender el cierre con la pregunta abierta;
% los párrafos finales del borrador sobre la reestructuración profunda
% de las relaciones de producción dentro de la ciencia conducen a ella.

¿Cómo puede ser apropiado este proceso de producción ultra socializado
y abstracto? La pregunta se deja abierta al lector --- y a la próxima
generación, que hereda no solo la pregunta de la apropiación, sino una
nueva conciencia y una nueva contradicción para pensar.

\begin{thebibliography}{9}

\bibitem{shaikh2016} Anwar Shaikh, \textit{Capitalism: Competition,
Conflict, Crises}, Oxford University Press, 2016.

\bibitem{marx1973} Karl Marx, \textit{Grundrisse: Foundations of the
Critique of Political Economy}, Penguin, 1973. (Fragmento sobre las
máquinas.)

\bibitem{marx1867} Karl Marx, \textit{Capital}, Vol.~1, 1867. Cap.~10,
La jornada de trabajo (la met\'afora del vampiro); Cap.~15, Maquinaria
e industria moderna.

\bibitem{tronti1966} Mario Tronti, \textit{Operai e capitale}, Einaudi,
1966. (La tesis f\'abrica--sociedad; la f\'abrica social.)

\bibitem{negri1979} Antonio Negri, \textit{Marx Beyond Marx: Lessons on
the Grundrisse}, 1979. Traducci\'on inglesa, 1984.

\bibitem{dallacosta1972} Mariarosa Dalla Costa y Selma James,
\textit{The Power of Women and the Subversion of the Community}, 1972.

\bibitem{federici2004} Silvia Federici, \textit{Caliban and the Witch:
Women, the Body and Primitive Accumulation}, Autonomedia, 2004.

\bibitem{debord1967} Guy Debord, \textit{La sociedad del espect\'aculo},
Buchet-Chastel, 1967.

\bibitem{lazzarato1996} Maurizio Lazzarato, ``Immaterial Labor,'' 1996,
en \textit{Radical Thought in Italy}, eds.\ Virno y Hardt, University
of Minnesota Press; ``From Capital-Labour to Capital-Life,''
\textit{Ephemera} 4(3), 2004.

\bibitem{althusser1970} Louis Althusser, ``Ideolog\'ia y aparatos
ideol\'ogicos de Estado,'' 1970, en \textit{Lenin and Philosophy and
Other Essays}, 1971.

\bibitem{baudrillard1981} Jean Baudrillard, \textit{Simulacros y
simulaci\'on}, Galil\'ee, 1981. Traducci\'on inglesa, University of
Michigan Press, 1994.

\bibitem{farjoun1983} Emmanuel Farjoun and Mosh\'e Machover,
\textit{Laws of Chaos: A Probabilistic Approach to Political Economy},
Verso, 1983.
% TODO: verificar la referencia a Farjoun, Machover y Zachariah citada
% en el borrador.

\bibitem{tao} Terence Tao, ``What does the advent of powerful AI models
mean for mathematicians like me?'' (ensayo).
% TODO: completar la referencia a Ghrist citada en el borrador.

\end{thebibliography}

\end{document}
